\documentclass[11pt,a4paper]{article}
\usepackage[margin=25mm]{geometry}
\usepackage{amsmath,amssymb,amsthm,booktabs}
\newtheorem{lemma}{Lemma}\newtheorem{theorem}{Theorem}
\title{Two maximal overgroups and the order 20 column}
\author{Human Compression project}\date{}
\begin{document}\maketitle
Let $G=O'N$, of order $460815505920$. We recover the corrected order 20 column
of the fixed Connor--Leemans Table~1. For a finite group $J$, write
\[
 x_J(p,q)=|\{(r,s)\in J^2:|r|=p,\ |s|=q,\ |rs|=2\}|,
\]
and let $e_J(p,q)$ count those pairs that generate $J$.
The complete subgroup roster, the named ATLAS maximal subgroup
structures and the previously acquired character evaluations are
explicit borrowed premises. Row numbers refer to the roster.
We do not infer an abstract group from an ambiguous extension label.

The argument has two parts. Direct incidence counts identify the two
maximal overgroups of every element of order 20. Their intersection
has a normal alternating subgroup with quaternion quotient; this
forces all pairs of the required kind into a direct product
$A_6\times C_4$. The last calculation uses the previously retained
360 permutations of $A_6$.

\begin{lemma}\label{two}
Every element $s$ of order 20 belongs to exactly two maximal subgroups
$A,B$ of $G$, in rows 7 and 11. Their intersection has row 32 and
order 2880. Conversely, each row 32 subgroup determines one such
pair $(A,B)$.
\end{lemma}
\begin{proof}
There is one conjugacy class of cyclic subgroups of order 20, row 493.
Following the complete immediate-overgroup lists gives exactly the
ancestor rows
\[
\begin{gathered}
1,7,8,11,17,30,32,33,34,47,57,58,59,101,133,143,164,\
185,208,209,210,211,235,312,403,493.
\end{gathered}
\]
Only rows 7 and 11 are maximal in $G$. This supplies existence of
both overgroups and of a row 32 subgroup containing $s$.

The row 7 subgroups are the centralizers of involutions. The accepted
centralizer calculation gives $|C_G(s)|=20$ for both element classes
of order 20. The cyclic subgroup $\langle s^5\rangle$ is a central
Sylow 2-subgroup of $C_G(s)$, hence its unique Sylow 2-subgroup.
Consequently $C_G(s)$ has just one involution, $s^{10}$, and a row 7
overgroup of $s$ must equal $C_G(s^{10})$.

There is one class of cyclic subgroups of order 5, row 574, of length
$640021536$. The row 11 maximal subgroups have class length
$17778376$. Their acquired trivial character counts 144 elements
of order 5, hence 36 cyclic subgroups of order 5. Counting incidences
in these two transitive conjugacy classes gives
\[
                    17778376\cdot36=640021536.
\]
Each cyclic subgroup of order 5 therefore lies in exactly one row 11
maximal subgroup. Apply this to $\langle s^4\rangle$.

A row 32 subgroup containing $s$ lies in one row 7 and one row 11
maximal subgroup, so these must be $A$ and $B$. It is maximal in $A$,
and $A\ne B$, whence it equals $A\cap B$. Finally the direct incidences
\[
              2857239\cdot56=17778376\cdot9=160005384
\]
show that every row 32 subgroup belongs to precisely one maximal
subgroup in each of rows 7 and 11. Here 160005384 is the row 32 class
length. These are single incidence counts; products along paths in
the subgroup diagram have not been used as containment numbers.
\end{proof}

\begin{lemma}\label{intersection}
Let $H$ be a row 32 subgroup. It has a normal subgroup $N\cong A_6$
with $H/N\cong Q_8$. Exactly one of its three subgroups of index two,
say $D$, contains elements of order 20, and $D\cong A_6\times C_4$.
If $s\in H$ has order 20 and $rs$ is an involution, then $r,s\in D$.
\end{lemma}
\begin{proof}
Let $B$ be the row 11 overgroup of $H$. The primary ATLAS entry
$(3^2:4\times A_6).2$ supplies a normal subgroup
\[
 B_0=F\times A_6,\qquad F=C_3^2:C_4,
\]
of index two in $B$. This named structural premise is separate from
the roster's notation for $H$. Since 2880 does not divide
$|B_0|=12960$, the subgroup $D=H\cap B_0$ has order 1440 and index
nine in $B_0$.

The projection of $D$ to $A_6$ has index dividing nine. Index three
is impossible, since a proper index-three subgroup would give a
nontrivial homomorphism from the simple group $A_6$ into $S_3$.
Index nine would give a subgroup of order 40. In such a subgroup the
number of Sylow 5-subgroups divides eight and is congruent to one
modulo five, so equals one. This would place the order 40 subgroup
inside $N_{A_6}(C_5)$, which has order ten. Indeed, the normalizer
of a 5-cycle subgroup in $S_6$ has order 20, and its automorphism
of exponent two is odd, leaving exactly half in $A_6$.
Thus $D$ projects onto $A_6$.

The group $F$ is solvable and $A_6$ is perfect. A sufficiently late
derived subgroup of $D$ has trivial projection to $F$ and full
projection to $A_6$, and therefore equals the factor $1\times A_6$.
It follows that $D=A_6\times K$ for an order four subgroup $K$ of
$F$. Every Sylow 2-subgroup of $F$ is conjugate to its cyclic
complement of order four, so $D\cong A_6\times C_4$.

The complete maximal-subgroup list for $H$ contains exactly three
index-two subgroups, one in each of rows 57, 60 and 61. Only row 57
is in the cyclic 20 ancestor list. Since $D$ contains elements of
order 20, it is that subgroup. Denote the other two by $E$ and $F_1$.
Each of $D,E,F_1$ has exactly one subgroup of index two, as the
complete roster lists show. In $D=A_6\times C_4$ this subgroup is
$K_0=A_6\times C_2$. The intersection of any two distinct index-two
subgroups of $H$ has index two in each. Uniqueness therefore gives
\[
             D\cap E=D\cap F_1=E\cap F_1=K_0.
\]
In particular $K_0$ is normal in $H$, and so is its derived subgroup
$N=K_0'\cong A_6$.

The quotient $H/N$ has order eight. Each of $D/N,E/N,F_1/N$ has
order four and a unique index-two subgroup, hence is cyclic.
Their pairwise intersections are $K_0/N$, of order two, so their
union has $2+3(4-2)=8$ elements and exhausts $H/N$. It follows that
all involutions of $H/N$ lie in the common subgroup $K_0/N$. Equivalently,
if $i,j$ generate two different cyclic subgroups of order four, then
$i^2=j^2=(ij)^2$ is the unique involution. These are the quaternion
relations, and $H/N\cong Q_8$.

For $|s|=20$, the image $sN$ has order four: if it had order at most
two, the subgroup $N\cong A_6$ would contain $s^2$, of order ten.
The other two index-two subgroups contain no element of order 20,
so $s\in D$. The image of an involution $rs$ in $Q_8$ is either the
identity or its central involution. Both belong to $D/N$, and
$rN=(rs)N(sN)^{-1}$ therefore belongs to $D/N$. Thus $r\in D$.
\end{proof}

\begin{lemma}\label{literal}
The only nonzero values of $x_H(p,20)$ are
\[
           x_H(4,20)=8640,\quad x_H(12,20)=5760,
           \quad x_H(20,20)=11808.
\]
\end{lemma}
\begin{proof}
By Lemma~\ref{intersection} we may count inside $D=A_6\times C_4$.
Write $s=(b,z^\varepsilon)$, with $|b|=5$ and
$\varepsilon\in\{1,3\}$. For $rs$ to have order two, the $C_4$
coordinate of $r=(a,z^\delta)$ must also be odd. If $ab$ is an
involution, all four choices of $(\delta,\varepsilon)$ work.
If $ab=1$, exactly two give a nonidentity involution and the other
two give the identity. Thus
\[
 x_H(p,20)=4\!\sum_{\operatorname{lcm}(j,4)=p}x_{A_6}(j,5)
               +288\,\mathbf1_{p=20}.
\]
The retained literal permutation set for $A_6$ has 144 elements of
order five. Its complete pair calculation gives
\[
\begin{array}{c|rrrr}
j&2&3&4&5\\\hline
x_{A_6}(j,5)&720&1440&1440&2880.
\end{array}
\]
For clarity about verification, the reader independently generates
all 360 elements using adjacent 3-cycles, checks that each is even,
and reconstructs every partner pair with the opposite composition
convention. It also enumerates all pairs in the 1440-element direct
product, including the cases with identity $ab$. These checks use
only permutation multiplication and exact element orders. The
stated values follow, and no other order $p$ occurs.
\end{proof}

\begin{theorem}
For every $p$ in the fixed table, one has
\[
 e_G(p,20)=x_G(p,20)-\sum_{M}[G:M]x_M(p,20)
                       +160005384\,x_H(p,20),
\]
where $M$ ranges over representatives of the thirteen maximal
subgroup classes. The resulting ordered automorphism-orbit counts
supply all sixteen entries of the corrected order 20 column.
\end{theorem}
\begin{proof}
A nongenerating pair is contained in a maximal subgroup. By
Lemma~\ref{two}, a pair with second element of order 20 is contained
in at most the two specified maximal subgroups. The pair is counted
twice in the maximal sum exactly when both entries lie in their
intersection $H$. Conversely, each row 32 subgroup determines
exactly one pair of maximal overgroups. Inclusion-exclusion therefore
gives the displayed identity.

The action of $\operatorname{Aut}(G)$ on generating pairs is free:
an automorphism fixing both entries fixes their generated group
pointwise. Its order is $2|G|$, by the already bound ATLAS premise.
Dividing the displayed identity by this order gives the following
values. The already accepted correction at $(20,31)$ is retained:
the exact value is 9118, whereas the source prints 7246. It is not a
new correction from this component. The values count ordered pairs; no extra division by duality is
made. The case displayed as $(20,31)$ uses the equivariant reversal
of the two entries, which preserves the order of their product.
\[\begin{array}{c|rr}
p& e_G(p,20)/|\operatorname{Aut}(G)|&\text{source}\\\hline
3&44&44\\
4&554&554\\
5&800&800\\
6&1948&1948\\
7&2916&2916\\
8&8880&8880\\
10&7224&7224\\
11&13094&13094\\
12&3984&3984\\
14&5024&5024\\
15&6424&6424\\
16&35644&35644\\
19&22572&22572\\
20&14238&14238\\
28&10292&10292\\
31&9118&7246\\
\end{array}\]

\end{proof}

\paragraph{Attribution and scope.}
The numerical claims originate in Connor and Leemans' fixed Table~1.
The complete subgroup roster, ATLAS normal subgroup structure, and
character tables remain borrowed inputs; the finite readers do not
prove those classifications. The two incidence arguments, the
quaternion quotient argument, and their application to the common
pair-counting identity are the present proof. The centralizer
criterion and the literal $A_6$ data were established in earlier
accepted components and are reused here. No priority claim is made
for the elementary group-theoretic arguments. Relative to the
accepted 92 entries, the new entries are $(4,20)$, $(12,20)$ and
$(20,20)$. This result supplies no new chirality assertion, Lean
certificate, or whole-original completion.
\end{document}
